\documentclass{ctexbook}
\usepackage{imakeidx}
\makeindex[title=索引,intoc]
\usepackage[]{subcaption}
\usepackage[]{bicaption}
\usepackage{graphicx}
\usepackage[]{amsmath}
\usepackage[]{amssymb}
\usepackage[]{hyperref}
\hypersetup{colorlinks, linkcolor=blue}
\usepackage[]{bigstrut}
\usepackage[dvipsnames,svgnames]{xcolor}
\usepackage[]{multicol}
% \usepackage{chemfig}
\usepackage{mhchem}
\usepackage[tone,extra,safe]{tipa}
\usepackage[normalem]{ulem}
\usepackage{tabularray}
\UseTblrLibrary{booktabs}
% \usepackage[]{paracol}
\usepackage[]{lipsum}
\usepackage[]{zhlipsum}
\usepackage[]{sidenotes}
\usepackage[]{subfiles}
\usepackage[]{siunitx}
\usepackage[]{metalogo}
\usepackage[]{tcolorbox}
\tcbuselibrary{skins}
\tcbuselibrary{breakable}
\tcbuselibrary{documentation}
\usepackage{geometry} 
\geometry{left=2.5cm,right=2.5cm,top=2.5cm,bottom=2.5cm}

\newtcolorbox[auto counter,number within=section]{sample}[2][]{%
arc is angular,breakable,parbox=false,before upper=\par,before lower=\par,drop shadow,skin=bicolor,colbacklower=white,colback=blue!8!white,colframe=blue!46!black, arc=1mm,fonttitle=\bfseries, title={\heiti Sample}~\thetcbcounter: #2,#1}

\newtcolorbox[auto counter,number within=section]{note}[2][]{%
breakable,parbox=false,before upper=\par,before lower=\par,drop shadow,skin={spartan,bicolor},colbacklower=white,colback=red!3!white,colframe=brown!46!black,fonttitle=\bfseries, title={\heiti Note}~\thetcbcounter: #2,#1}

\newtcolorbox[auto counter,number within=section]{task}[2][]{%
colback=Emerald!10,colframe=cyan!40!black,fonttitle=\bfseries,title={\heiti Task}~\thetcbcounter: #2,#1}

  
\usepackage[]{adforn}
% \renewcommand{\chaptername}{第 \thechapter{} 章}
% \newcommand{\exercisename}{Exercice}
\newenvironment{problemset}[1][\chaptername~\;Exercice]{
  \begin{center}
    \phantomsection\addcontentsline{toc}{section}{\texorpdfstring{chapter\;Exercice}{Exercice}}
    % \markboth{#1}{\rightmark}
    \markright{#1}
    \textcolor{black}{\Large\bfseries\adftripleflourishleft~#1~\adftripleflourishright}
  \end{center}
  \begin{enumerate}}{
  \end{enumerate}}


\begin{document}


\begin{titlepage}
    \vspace*{\stretch{1}}
    \begin{center}
      {\huge\bfseries Abstract Writing for Medical \\\vspace{3pt} Research Papers}\\[3ex]
      {\huge\bfseries 医学论著英语摘要写作}\\[6.5ex]
      {\large\bfseries Qi Hui, Chen Feina, Guo Haiyan}           \\
      \vspace{4ex}
    %   Thesis  submitted to                    \\[5pt]
      \textit{Fujian Medical University}                \\[2cm]

    %   \includegraphics[width=0.4\textwidth]{fig/fjmulogo.jpg}

    %   in partial fulfilment for the award of the degree of \\[2cm]
    %   \textsc{\Large Doctor of Philosophy}    \\[2ex]
    %   \textsc{\large Mathematics}             \\[12ex]
      \vfill
      Department of Arts and Sciences\\
    %   Address                                 \\
      \vfill
      \date{}
    \end{center}
    \vspace{\stretch{2}}
\end{titlepage}

\frontmatter
\chapter*{译者序}

本书是由福建医科大学文理艺术学院的齐晖、陈菲娜、郭海燕老师为主编，陈晶为学术秘书，交由复旦大学出版社出版，专供福医大学生使用的英语摘要写作教科书。我本人也是三位老师的学生，纵然课堂生动有趣、干货满满，但苦于同校前辈制作的扫描件观感不佳，笔记整理不便，译者决心要进行文字重排处理。其中自觉原书排版不善之处，皆进行重新编排，以符合译者审美。

原书通本以英文编写，编者似乎意图借此提升我等英语阅读水平，奈何文本中穿插语言学专有名词，初学时疲于翻译、苦不堪言。此外，期末复习期间，全英文本并不利于提升复习效率，故译者对主要文本进行翻译，对照复习。本套重置本将基于该译本进行整理，包含三种排版样式——原文重排版、双语对照版、译文版。以上三个版本请学弟学妹们按需取用，20届学长祝各位期末考试顺利。

本书为个人翻译作品，若发现纰漏，请在GitHub上提交\verb|issue|。在此声明，此书仅供学习交流使用，请勿用于商业用途或其他领域。

\vfill

\begin{flushright}
    neurocane\\
    2023/7/7
\end{flushright}

\newpage

\section*{编译环境}

\begin{itemize}
    \item \textbf{操作系统}： Windows
    \item \textbf{语言}： \LaTeX
    \item \textbf{编译环境}： \XeLaTeX
    \item \textbf{TeX Live版本}： TeX Live 2022
\end{itemize}

\tableofcontents

\mainmatter

\newgeometry{includeall,left=2cm,right=1cm,top=2.5cm,bottom=2.5cm,textwidth=12cm,marginparwidth=5cm}

\chapter{Overview of Abstracts}\label{chapter1}
\section{Definition of an Abstract}

The American Psychological Association (APA) Style (2010) states that an abstract is a brief, comprehensive summary of the content of an article. According to the American National Standards Institute (1979), an abstract is an abbreviated accurate representation of the content of a document, preferably prepared by its author(s) for publication with it. In general, an abstract is a concise, accurate and comprehensive statement of the content of an article. It is original rather than excerpted.

\section{Importance and Functions of an Abstract}

An abstract is a distinct genre, and to some extent plays a pivotal role in academic readingand writing. In the era of information explosion, an enormous number of new publications are produced in the academic community each day. There is no practical way for every reader to get access to every new article, or to read every new publication even if it is accessible. Theabstracts published online, which are concise and comprehensive, can be obtained easily andquickly. Abstract reading, then, may be a useful starting point of any academic reading and writing. In this sense, an abstract is the most read part of an article.\sidenote{劳仑衣普桑，认至将指点效则机，最你更枝。想极整月正进好志次回总般，段然取向使张规军证回，世市总李率英茄持伴。}

An abstract has at least three functions (Huckin, 2001). First, it serves as a stand-alone mini text, giving readers a quick summary of a study's objectives, methodology, findings and conclusions, which are the major components of abstracts. Second, it serves as a screening device, and gives readers an adequate view on whether the full-length article is of great value to their needs and worth further reading. A good abstract, to some extent, increases the chance of being cited or referenced. Third, for those readers who do opt to read the article as a whole, the abetract serves as a preview, creating an interpretive frame that can guide reading.

\section{Types of Abstracts}

Generally, abstracts fall into two categories, indicative and informative, depending on the type of information they convey. A typical distinction between them is that the indicative abstract, viewed as the outline of the paper, is usually shorter and simpler, while the informative abstract, viewed as the summary of the paper is usually longer and more thoroagh. 

These two types of abstracts also differ in the components they contain. Indicative abstracts often include the purpose, scope, and methods of the report or study, but seldom include theresults or conclusions. Reading indicative abstracts could not substitute reading the paper, because not all the crucial components are covered. It is more widely used in social science papers. On the other hand, informative abstracts usually include all the crucial components of the study, such as the background, purpose, methods, results, and conclusions. It is the type of abstracts widely used in medical field. In this book, we focus on the writing of informative abstracts. The abstracts referred to in the following chapters are informative abstracts.

\section{Types of Informative Abstracts}

There are two types of informative abstracts, structured abstracts and unstructured abstracts.

Where a heading or label is used at the beginning of the text in each section, it is a structured abstract. Each section is usually written in a separate paragraph, but sometimes sections are written in a sole or continuous paragraph. Headings might be background, objectives, methods, results, conclusions, and so on. They vary according to the criteria set by different journals. Structured abstracts appear to be favored by medically-relevant publications.

Where no heading or label is used to indicate different parts of an abstract, it is an unstructured abstract. It is always a sole paragraph. The major difference between the two types of abstracts lies in whether there are headings or not. In an unstructured abstract the content and sequence of the items are written as it is in the structured one.

Journals mandate which style should be used, so check the author guidelines if you' re not sure. If it is not mentioned, keep an eye out for the type of abstracts preferable in the journals where you are willing to have your paper submitted and published. Write your abstracts in the style which dominates.

% \tcbset{colback=blue!8!white,colframe=blue!46!black,
% arc=1mm}
% \begin{tcolorbox}[arc is angular,skin=bicolor,colbacklower=white,title={\heiti 例题}]
% 这是一个 \textbf{tcolorbox}。
% \tcblower
% 这个是下部
% \end{tcolorbox}


\begin{sample}[label={myautocounter}]{\heiti}
  \textbf{BACKGROUND} 
  
  In patients with acute heart failure, early intervention with an intravenous vasodilator has been proposed as a therapeutic goal to reduce cardiac-wall stress and, potentially, myocardialinjury, thereby favorably affecting patients' long-term prognosis.

  \textbf{METHODS} 
  
  In this double-blind trial, we randomly assigned 2, 157 patients with acute heart failure to receive a continuous intravenous infusion of either ularitide at a dose of 15 ng per kilogram of body weight per minute or matching placebo for 48 hours, in addition to accepted therapy.
  Treatment was initiated a median of 6 hours after the initial clinical evaluation. The coprimaryd outcomes were death from cardiovascular causes during a median follow-up of 15 months and a hierarchical composite end point that evaluated the initial 48-hour clinical course.
  
  \textbf{RESULTS} 
  
  Death from cardiovascular causes occurred in 236 patients in the ularitide group and 225 patients in the placebo group (21.7\% vs. 21.0\%; hazard ratio, 1.03;96\% confidence interval, 0.85 to 1.25; P=0.75). In the intention-to-treat analysis, there was no significant between-group difference with respect to the hierarchical composite outcome. The ularitide group had greater reductions in systolic blood pressure and in levels of N-terminal pro-brain natriuretic peptide than the placebo group. However, changes in cardiac troponin T levels during the infusion did not differ between the two groups in the 55\% of patients with paired data.

  \textbf{CONCLUSIONS} 
  
  In patients with acute heart failure, ularitide exerted favorable physiological effects (without affecting cardiac troponin levels), but short-term treatment did not affect a clinical composite end point or reduce long-term cardiovascular mortality.

  \begin{flushright}
    ---Effect of Ularitide on Cardiovascular Mortality in Acute Heart Failure.

    \emph{New England Journal of Medicine (2017)}
  \end{flushright}

  
\end{sample}

\begin{sample}[label={myautocounter}]{\heiti}
  \textbf{OBJECTIVE} 
  
  To evaluate the association between the parameters of 24-hour multichannel intraluminal impedance (MII)-pH monitoring and the symptoms or quality of life (QoL) in laryngopharyngeal reflux (LPR) patients.

  \textbf{DESIGN} 
  
  Prospective cohort study without controls.

  \textbf{SETTING} 
  
  University teaching hospital.

  \textbf{METHODS} 
  
  Forty-five LPR patients were selected from subjects who underwent 24-hour MII-pH monitoring and were diagnosed with LPR from September 2014 to May 2015. Reflux Symptom Index (RSI), Health-related Quality of Life (HRQoL), Short Form 12 (SF-12) Survey questionnaires were surveyed. Spearman's correlation was used to analyse the association between the symptoms or QoL and 24-hour MII-pH monitoring.

  \textbf{RESULTS}

  Most parameters in 24-hour MII-pH monitoring showed weak or no correlation with RSI, HRQoL and SF-12. Only number of non-acid reflux events that reached the larynx and pharynx (LPR-non-acid) and number of total reflux events that reached the larynx and pharynx (LPR-total) parameters showed strong correlation with heartburn in RSI (R=0.520, $P<0.001$, R=0.478, $P=0.001$, respectively). Multiple regression analysis showed that there was only one significant regression coefficient between LPR-non-acid and voice/hoarseness portion of HRQoL (b=1.719, $P=0.022$).

  \textbf{CONCLUSION} 

  Most parameters of 24-hour MII-pH monitoring did not reflect subjective symptoms or QoL in patients with LPR.

  \begin{flushright}
    ---Association between 24-hour combined multichannel intraluminal impedance-pH monitoring and symptoms or quality of life in patients with laryngopharyngeal reflux. 
    
    \emph{Clinical Otolaryngology (2017)}
  \end{flushright}

  
\end{sample}

\begin{sample}[label={myautocounter}]{\heiti}
  Due to the high incidence of recurrent squamous cell carcinoma of the head and neck andd the toxicity profile of current salvage regimens, there is a need for tolerable and effective treatment options. We performed a retrospective matched case series to report our experience with recurrent high-risk patients who received capecitabine (CAP) therapy in the adjuvant setting after salvage therapy. The 5-year recurrence-free survival rates for the CAP and control cohorts were 54\% (95\% CI, 0.27\%--0.75\%) and 27\% (95\% CI, 0.09\%--0.50\%), respectively.Multivariable Cox modeling showed a significant improvement in recurrence-free survival in thed CAP cohort (hazard ratio, 0.19; 95\% CI, 0.04--0.92; $P=.0 392$). While this was a respective analysis that could not control for all variables, these exploratory findings offer insights that may inform a prospective study to determine CAP efficacy.
  
  \begin{flushright}
    ---Capecitabine after Surgical Salvage in Recurrent Squamous Cell Carcinoma of Head and Neck. 
    
    \emph{Otolaryngology---Head \& Neck Surgery (2017)}
  \end{flushright}


\end{sample}

\begin{note}[label={myautocounter}]{\heiti Corpus used for this book}
  The data used and analyzed in this book are from a custom-built corpus with 1.15 million tokens of medical research article (RA) abstracts. The discipline of medicine is divided into 18 sub-disciplines, and RA abstracts from 2 to 3 leading journals are randomly retrieved in each sub-discipline with relatively similar number of texts for each sub-discipline (\autoref{tab:Sub-disciplines and journals in each sub-discipline}). The journals selected are all with relatively high impact factors.

\end{note}

{\small
  \begin{longtblr}[
      caption = {Sub-disciplines and journals in each sub-discipline},
      label = {tab:Sub-disciplines and journals in each sub-discipline},
  ]{
      width = \textwidth,
      colspec = {X[1,l,h]  X[2,l,h]},
      rowhead = 1, rowfoot = 0, % 每个分页里表头表尾的数量
      % row{odd} = {blue8}, 
      row{even} = {LemonChiffon},
      column{2} = {font=\itshape}
  }
    
    \toprule
    Sub-discipline & Journal \\
    \midrule

    Anesthesiology & {British Journal of Anaesthesia \\ Anesthesiology \\ Anesthesia and Analgesia}\\
    Dermatology & {Journal of American Academy of Dermatology \\ Giornale Italiano di Dermatologia e Venereologia}\\
    Emergency Medicine & {Annals of Emergency Medicine \\ Internal and Emergency Medicine \\ Academic Emergency Medicine}\\
    Geriatrics & {Neurobiology of Aging \\ Aging Cell \\ Age and Ageing}\\
    Internal Medicine & {The New England Journal of Medicine \\ The Lancet \\ JAMA-Journal of the American Medical Association}\\
    Medical Imaging & {The Journal of Nuclear Medicine \\ Investigative Radiology \\ Radiology}\\
    Medical Laboratory & {Clinical Chemistry and Laboratory Medicine \\ Clinical Biochemistry}\\
    Neurology & {The Lancet Neurology \\ Annals of Neurology}\\
    Obstetrics and Gynecology & {Obstetrics \& Gynecology \\ American Journal of Obstetrics \& Gynecology \\ An International Journal of Obstetrics \& Gynecology}\\
    Oncology & {Journal of Clinical Oncology \\ The lancet Oncology}\\
    Ophthalmology & {Ophthalmology \\ American Journal of Ophthalmology \\ Archives of Ophthalmology}\\
    Otolaryngology (ENT) & {Head \& Neck \\ Clinical Otolaryngology \\ Otolaryngology---Head \& Neck Surgery}\\
    Pain Medicine & {The Clinical Journal of Pain \\ Pain Medicine \\ Regional Anesthesia and Pain Medicine}\\
    Pediatrics & {Journal of the American academy of child \& Adolescent psychiatry \\ Pediatrics \\ JAMA pediatrics}\\
    Physical medicine and rehabilitation & {Neurorehabilitation and neural repair \\ Journal of fluency disorders}\\
    Psychiatry & {Molecular psychiatry \\ The American journal of psychiatry \\ JAMA psychiatry}\\
    Sports medicine & {Medicine and Science in Sports and Exercise \\ Sports Medicine \\ The American Journal of Sports Medicine}\\
    Surgery & {Annals of Surgery \\ American Journal of Transplantation \\ Journal of Neurology, Neurosurgery \& Psychiatry}\\

    \bottomrule

  \end{longtblr}
  }

  \begin{task}[label={myautocounter}]{\heiti Corpus-based task}
    Can you build your own corpus with at least 100, 000 tokens?
  \end{task}

  \begin{note}[label={myautocounter}]{\heiti Corpus used for this book}
    In this book, "corpus-based tasks" are designed to enhance your ability to explore language realizations of medical RA abstracts with corpus approach. Most of these tasks might require the use of software such as AntConc, WordSmith, and so on.
  \end{note}

  \section{Glossary}

  {\small
  \begin{longtblr}[
      caption = {Glossary of Chapter 1},
      label = {tab:Glossary of Chapter 1},
  ]{
      width = \textwidth,
      colspec = {X[1,l,h]  X[1,l,h]  X[3,l,h]},
      rowhead = 1, rowfoot = 0, % 每个分页里表头表尾的数量
      % row{odd} = {blue8}, 
      row{even} = {azure9},
  }
      
  \toprule
  \textbf{WORDS} & \textbf{MEANING} & \textbf{MEANING OR EXAMPLE}\\
  \midrule

  \textbf{excerpt}/\textipa{"eks3:pt}/ & \emph{v.} 摘录；引用  & If a long piece of writing or music is excerpted, short pieces from it are printed or played on their own. \\
  \textbf{genre}/\textipa{"ZA:nr@}/ & \emph{n.} 体裁 & a particular type of art, writing, music etc, which has certain features that all examples of this type share.\\
  {\textbf{mandate} \\ /\textipa{"m\ae ndeIt}/} & \emph{v.} 授权；强制执行；委托办理 & to tell someone that they must do a particular thing. \\
  {\textbf{methodology} \\ /\textipa{""meT@"d6l@\textdyoghlig i}/} & \emph{n.} 方法学 & a set of methods and principles used to perform a particular activity. \\
  \textbf{opt}/\textipa{6pt}/ & \emph{v.} 选择；挑选 & to choose one thing or do one thing instead of another \\
  \textbf{pivotal}/\textipa{"pIv@tl}/ & \emph{adj.} 关键性的；核心的 & more important than anything else in a situation or system. \\
  
  \bottomrule

  \end{longtblr}
  }

  \begin{problemset}
    \item \textbf{Identify whether the following abstracts are structured or unstructured and tell the reasons.}
    
    \textbf{Abstract 1}

    \hspace*{2em}Objectives: The aim of this study was to analyze changes in health care utilization and cost among a sample of highly impaired children and adolescents who sought a 3-week intensive interdisciplinary pain treatment (IPT).

    \hspace*{2em}Materials and Methods: Claims data from 7 statutory health insurance companies were analyzed for 65 children and adolescents who sought IIPT at the German Paediatric Pain Centre. The annual health care utilization and cost were determined for the following 4 areas: outpatient care, inpatient care, medications, and remedies and aids. We analyzed the changes in resource utilization in the year before (pre\_1 y) IPT and in the subsequent year (post\_1 y).
    
    \hspace*{2em}Results: Within the first year after IPT, overall health care costs did not decrease significantly. However, the pattern of health care utilization changed. First, significantly more children and adolescents started outpatient psychotherapy ($P=0.001$). Second, the number of hospitalized children decreased significantly from 1-year pre to 1-year post ($P=0.001$). Accordingly, there were significantly fewer hospitalizations for primary chronic pain disorders at 1-year post ($P<0.001$). The prescription of nonopioids, co-analgesics and opioids was significantly reduced from 1-year pre to 1-year post (all $P<0. 013$).
    
    \hspace*{2em}Discussion: The present results indicate that the health care costs of children and adolescents with severe chronic pain disorders do not significantly decrease 1 year after IPT; however, the treatment becomes more goal-focused. Differential diagnosis measures and nonindicated therapeutic interventions decreased, and more indicated interventions, such as psychotherapy, were used. Future research is needed to investigate the economic long-term changes after IPT.
    
    \begin{flushright}
      ---Health Care Utilization and Cost in Children and Adolescents with Chronic Pain: Analysis of Health Care Claims Data 1 Year Before and After Intensive Interdisciplinary Pain Treatment. 
      
      \emph{The Clinical Journal of Pain (2017)}
    \end{flushright}

    \textbf{Abstract 2}

    \hspace*{2em}Previous studies of brain structure in Tourette syndrome (TS) have produced mixed results, and most had modest sample sizes. In the present multicenter study, we used structural magnetic resonance imaging (MRI) to compare 103 children and adolescents with TS to a well-matched group of 103 children without tics. We applied voxel-based morphometry methods to test gray matter (GM) and white matter (WM) volume differences between diagnostic groups, accounting for MRI scanner and sequence, age, sex and total GM$+$WM volume. The TS group demonstrated lower WM volume bilaterally in orbital and medial prefrontal cortex, and greater GM volume in posterior thalamus, hypothalamus and midbrain. These results demonstrate evidence for abnormal brain structure in children and youth with TS, consistent with and extending previous findings, and they point to new target regions and avenues of study in TS. For example, as orbital cortex is reciprocally connected with hypothalamus, structural abnormalities in these regions may relate to abnormal decision making, reinforcement learning or somatic processing in TS.

    \begin{flushright}
      ---Brain structure in pediatric Tourette syndrome. 

      \emph{Molecular Psychiatry (2017)}
    \end{flushright}

    \textbf{Abstract 3}

    \hspace*{2em}Objective: To assess the feasibility of detecting signature volatile organic compounds in the breath of patients with oral squamous cell carcinoma.
    
    \hspace*{2em}Study Design: Prospective cohort pilot study.
    
    \hspace*{2em}Setting: University hospital.
    
    \hspace*{2em}Subjects and Methods: Using gas chromatography and mass spectrometry, emitted volatile organic compounds in the breath of patients before and after curative surgery (n=10) were compared with those of healthy subjects (n=4). It was hypothesized that certain volatile organic compounds disappear after surgical therapy. A characteristic signature of these compounds for diseased patients was compiled and validated.
    
    \hspace*{2em}Results: Breath analyses revealed 125 volatile organic compounds in patients with oral cancer. A signature of 8 compounds that were characteristic for patients with oral cancer could be detected: 3 from this group presented were absent after surgery.
    
    \hspace*{2em}Conclusion: The presented results confirmed the hypothesis of an absence of cancer-associated volatile organic compounds in the breath after therapy. In this pilot study, we proved the feasibility of this test approach. Further studies should be initiated to establish protocols for usage in a clinical setting.
    
    \begin{flushright}
      ---Volatile Organic Compounds in the Breath of Oral Squamous Cell Carcinoma Patients: A Pilot Study. 

      \emph{Otolaryngology Head and Neck Surgery (2017)}
    \end{flushright}
   
  \end{problemset}


\chapter{Move and Step Identification}\label{chapter2}

\section{Move Identification}

A medical RA abstract consists of moves which work together to achieve its communicative purposes. A move in this sense is ``a section of a text that performs a specific communicative function'' (Kanoksilapatham, 2007, p.23).

A four move scheme is used in most medical RA abstracts, with move 1 (M1) creating a research space move 2 (M2) describing research process, move 3 (M3) summarizing principal results and move 4 (M4) drawing conclusions. All of them are conventional moves in medical RA abstracts.

In a structured abstract, M2, M3 and M4 could be easily recognized via the headings of methods, results and conclusions. Although there is no specific heading of ``methods'' in some abstracts, M2 is subdivided into several steps which could be clearly recognized. They could be labelled ``design''; ``setting'', ``participants'', `` interventions'', `` main outcome measures", and so on.

When it comes to M1, sections labelled with either ``objectives" or ``background'' or both are included. In some structured abstracts, both ``background'' and ``"objectives"'' are labelled, in some only ``objectives", and in others only ``background". Those sections with the label of either ``objectives'' or ``background'' are usually comprised of both of them in terms of actual contents. Moreover, in practice, objectives and background are closely related and usually viewed as a whole to provide a specific communicative purpose of creating research space.

For a structured abstract, four moves could be recognized by headings, one of the lexical signals. For unstructured abstracts in which there are no headings or labels, manual recognition of four moves is needed. Other lexical signals could be helpful in recognizing the moves, which are illustrated in \hyperref[chapter3]{Chapter 3} to \hyperref[chapter6]{Chapter 6}.

\begin{sample}[label={myautocounter}]{\heiti}

\textbf{(M1)}

\textbf{BACKGROUND }

it is still equivocal whether there is a potential role of late-life physical activity in ameliorating the challenges of increasing healthcare expenditure due to the consequence of global population ageing.

\textbf{OBJECTIVE} 

this study aimed to examine the prospective association between physical activity and subsequent hospital care utilisation in older adults and to explore the optimal dose of physical activity required to reduce hospital care utilisation.

\textbf{(M2)}

\textbf{DESIGN} 

this was a prospective cohort study based on the data from the Taiwan 2005 National Health Interview Survey, which were linked to the 2005-12 claims data from the National Health Insurance system.

\textbf{PARTICIPANTS}

1,760 older adults aged 65 or more.

\textbf{METHODS} 

the frequency, duration and intensity for physical activity were assessed, and total physical activity energy expenditure was estimated. The average annualised hospital care utilisation for the period 2006 through 2012, including number of hospitalisations, number of days in hospital and the costs of hospitalisation, were calculated.

\textbf{(M3)}

\textbf{RESULTS }

older adults engaging in at least moderate volume of physical activity ($\geqslant 1,000$ kcal/week) experienced fewer subsequent hospital admissions and fewer days in hospital than did sedentary individuals, after adjusting for covariates. Trends for reduced hospitalisation costs were also found. These associations persisted in sensitivity analyses, including tests of reverse causation.

\textbf{(M4)}

\textbf{CONCLUSION}

this study has provided evidence that older adults who are at least moderately active may minimise utilisation of hospital care services. The findings highlight the importance of maintaining a physically active lifestyle in later life.

\begin{flushright}
  ---Prospective association between late-life physical activity and hospital care utilisation: a 7-year nationwide follow-up study. \emph{Age and Aging (2017)}
\end{flushright}

\end{sample}

\begin{sample}[label={myautocounter}]{\heiti}

\textbf{(M1)}

\textbf{PURPOSE }

Stuttering can trigger anxiety and other psychological and emotional reactions, and limit participation in society. It is possible that psychological counseling could enhance stuttering treatment outcomes; however, little is known about how clients view such counseling. The purpose of this study was to gain an understanding of clients' experiences with, and perceptions of, a psychological counseling service that was offered as an optional adjunct to speech therapy for stuttering.

\textbf{(M2)}

\textbf{METHOD }

Nine individuals who stutter (13--38 years old) participated in semi-structured interviews. Six participants had taken part in psychological counseling; three participants did not do so. Interview data were analyzed using grounded theory as a guiding framework.

\textbf{(M3)}

\textbf{RESULTS }

Four thematic clusters emerged from participants' accounts: insights into personal decision-making, why others may not participate in counseling, psychological counseling as a worthwhile part of therapy, and counseling as a necessary component in a stuttering treatment program.

\textbf{(M4)}

\textbf{CONCLUSION}

In addition to experiencing barriers and facilitators to help-seeking that are reported in related fields, participants accounts also revealed novel facilitators (i.e., a `why not' mentality and the importance of having a pre-existing relationship with the clinician who offered the service) and barriers (i. e., viewing the service as a `limited resource,' and, the overwhelming nature of intensive stuttering treatment programs). Findings suggest that clients value the option to access psychological counseling with trained mental health professionals to support the stuttering treatment provided by speech-language pathologists. Participants made recommendations for the integration of psychological counseling into stuttering treatment programs.

\begin{flushright}
  ---Psychological counseling as an adjunct to stuttering treatment: Clients' experiences and perceptions. 
  
  \emph{Journal of Fluency Disorder (2017)}
\end{flushright}


\end{sample}


\begin{sample}[label={myautocounter}]{\heiti}
\textbf{(M1)}

Genetic and neuroimaging research has identified neurobiological correlates of obesity. However, evidence for an integrated model of genetic risk and brain structural alterations in the pathophysiology of obesity is still absent.

\textbf{(M2)}

Here we investigated the relationship between polygenic risk for obesity, gray matter structure and body mass index (BMI) by the use of univariate and multivariate analyses in two large, independent cohorts ($n=330$ and $n=347$).

\textbf{(M3)}

Higher BMI and higher polygenic risk for obesity were significantly associated with medial prefrontal gray matter decrease, and prefrontal gray matter was further shown to significantly mediate the effect of polygenic risk for obesity on BMI in both samples.

\textbf{(M4)}

Building on this, the successful individualized prediction of BMI by means of multivariate pattern classification algorithms trained on whole-brain imaging data and external validations in the second cohort points to potential clinical applications of this imaging trait marker.

\begin{flushright}
  ---Prefrontal gray matter volume mediates genetic risks for obesity. \emph{Molecular Psychiatry (2017)}
\end{flushright}

\end{sample}

\section{Step Identification}

In the genre with an obvious hierarchical structure, moves are usually composed of the steps or sub-moves, which are the subordinate units. In medical RA abstracts, some steps are conventional steps and others are optional steps. The type and frequency of the steps also show the rhetorical purpose of the author. Lexical signals could be helpful in recognizing the steps, which are illustrated in \hyperref[chapter3]{Chapter 3} to \hyperref[chapter6]{Chapter 6}.

The abstracts of different journals in the corpus are randomly extracted, with 6 articles in each sub-discipline, totaling 108 articles. Through manual recognition, the steps and communication functions that constitute each step are established (Table 1). If the step is used in more than 80\% of articles, it is considered conventional, otherwise optional.


\chapter{}\label{chapter3}
\chapter{}\label{chapter4}
\chapter{}\label{chapter5}
\chapter{}\label{chapter6}

\restoregeometry

\backmatter



	
\end{document}